\documentclass[12pt]{article}

% revise margins
\setlength{\headheight}{0.0in}
\setlength{\topmargin}{0.0in}
\setlength{\headsep}{0.0in}
\setlength{\textheight}{9.0in}
\setlength{\footskip}{0.0in}
\setlength{\oddsidemargin}{0.0in}
\setlength{\evensidemargin}{0.0in}
\setlength{\textwidth}{6.5in}

\pagestyle{empty}
\usepackage{times}
\usepackage{amsmath}

\begin{document}

\noindent \textbf{Biostat 140.668} \hfill Problem set 5

\bigskip 

\noindent \textbf{QTL mapping: other single-QTL models}

\begin{enumerate}

\item Consider mapping a simple Mendelian disease gene using a
backcross.  Suppose you have essentially complete genotype data, so
that the exact locations of crossovers can be determined.  Let $L$
denote the length (in cM) of the interval to which you will have
mapped the gene with $n$ backcross individuals (i.e., the distance
between the two closest flanking recombination events).

\begin{enumerate}
\item How many backcross individuals will you need so that $\text{E}(L) \le 1
\text{ cM}$?

\item How many backcross individuals will you need so that $\text{Pr}(L \le 1
\text{ cM}) \ge 80\%$?
\end{enumerate}

\item Derive the extension of the Kruskal-Wallis statistic
discussed in lecture, for nonparametric interval mapping.  

Let $R_i$
denote the ranks of the phenotypes, and let $p_{ig}$ denote the
probability that mouse $i$ has QTL genotype $g$, given the available
marker data.  Define $\bar{S}_g = \sum_i R_i p_{ig} / \sum_i p_{ig}$

\begin{enumerate}
\item Show that $\text{E}(\bar{S}_g \; | \; \text{marker data, no linkage}) =
(n+1)/2$.

\item Show that 
$$\text{var}(\bar{S}_g \; | \; \text{marker data, no
linkage}) = \left[ \frac{n+1}{12} \right] \; \left[\frac{n \sum_i
p_{ig}^2 - (\sum_i p_{ig})^2}{(\sum_i p_{ig})^2} \right] $$

\item Put it all together.
\end{enumerate}


\end{enumerate}




\end{document}
